% !Mode:: "TeX:UTF-8"
% Recommended compilation methods:
% 1. pdfLaTeX (recommended)
% 2. XeLaTeX

\documentclass{apmcmthesis}

\usepackage{url}

%%%%%%%%%%%% Contest Information %%%%%%%%%%%%%%%%%%%%%%%%%%
\tihao{A}                            % Problem choice
\baominghao{apmcm253XXXXX}           % Team control number

\begin{document}

\pagestyle{frontmatterstyle}

\begin{abstract}
The \verb|apmcmthesis| \LaTeX{} template is provided by \url{https://www.latexstudio.net} for the Asia-Pacific Mathematical Contest in Modeling (\url{http://www.apmcm.org}). It allows teams to concentrate on content rather than formatting.

Basic \LaTeX{} knowledge is assumed (references, equations, figures, tables, etc.). All commonly used packages are pre-loaded.

The latest template is maintained at \url{https://github.com/latexstudio/APMCMThesis}.

\keywords{\LaTeX \quad Mathematical Modeling \quad APMCM \quad Thesis Template}
\end{abstract}

% If the contest requires a Chinese abstract, delete the following block or translate it separately.
% \begin{abstract*}
% Chinese abstract here (optional)
% \keywords{keywords in Chinese}
% \end{abstract*}

\newpage
\tableofcontents

\newpage
\pagestyle{mainmatterstyle}
\setcounter{page}{1}

\section{Introduction}
To clarify the origin of the problem, the following background information is provided.

\section{Problem Restatement}

\subsection{How do we approximate the entire process?}

\begin{itemize}
  \item Point 1
  \item Point 2
  \item Point 3
\end{itemize}

\subsection{How do we define the optimal configuration?}

\begin{enumerate}
  \item From the perspective of efficiency:
  \item From the perspective of cost:
  \item Trade-off between the two:
\end{enumerate}

\subsection{Local Optimization vs. Global Optimization}

\begin{itemize}
  \item Local optimization
  \item Global optimization
  \item Practical balance
\end{itemize}

\subsection{Differences in Weights and Sizes}

\subsection{What if no data are available?}

\section{Models}

\subsection{Basic Model}

\subsubsection{Notation and Definitions}
Most symbols and definitions are derived from queuing theory.

\subsubsection{Assumptions}

\begin{enumerate}
  \item Assumption 1
  \item Assumption 2
  \item Assumption 3
\end{enumerate}

\subsubsection{Model Formulation}

1) Utility function
\begin{itemize}
  \item Waiting cost:
  \item Congestion loss:
  \item Weight of each factor:
  \item Trade-off strategy:
\end{itemize}

% Remove or replace the original promotional images
%\begin{figure}[!ht]
%  \centering
%  \includegraphics[width=3cm]{example-image-a} \quad \includegraphics[width=5cm]{example-image-b}
%  \caption{Follow our updates for more resources (optional)}
%\end{figure}

3) Global vs. local optimization
\begin{itemize}
  \item Global optimization:
  \item Local optimization:
  \item Optimal number of servers:
\end{itemize}

\subsubsection{Solution and Results}

\begin{enumerate}
  \item Integer programming solution
  \item Numerical results
\end{enumerate}

\subsubsection{Result Analysis}

\begin{itemize}
  \item Comparison of local and global optimization
  \item Sensitivity analysis: results are sensitive to the three key parameters
  \item Parameter trends
  \item Comparison with existing designs
\end{itemize}

\subsubsection{Strengths and Weaknesses}

\begin{description}
  \item[Strengths:] The improved model addresses factors neglected in the basic model. Results show it is more reasonable and significantly outperforms current designs.
  \item[Weaknesses:] The model is still a large-scale approximation, limiting its use in high-precision scenarios.
\end{description}

\section{Conclusions}

\subsection{Summary of Results}

\begin{itemize}
  \item Key finding 1
  \item Key finding 2
  \item Key finding 3
\end{itemize}

\subsection{Methods Employed}

\begin{itemize}
  \item Queuing theory
  \item Integer programming
  \item Sensitivity analysis
\end{itemize}

\subsection{Applications}

\begin{itemize}
  \item Application 1
  \item Application 2
  \item Application 3
\end{itemize}

\section{Future Work}

\subsection{Alternative Approaches}

\subsubsection{Limitations of Queuing Theory}
\subsubsection{Other Potential Models}
\subsubsection{Data-Driven Methods}
\subsubsection{Dynamic Optimization}

\section*{References}
\begin{thebibliography}{9}
\bibitem{ref1} Author(s), Book/Journal Title, Publisher/Journal, Year.
\bibitem{latexstudio} \LaTeX{} Resources and Tutorials, \url{https://www.latexstudio.net}
\bibitem{github} APMCMThesis Template Repository, \url{https://github.com/latexstudio/APMCMThesis}
\end{thebibliography}

\newpage
\section*{Appendix}

\begin{lstlisting}[language=matlab,caption={MATLAB source code of the algorithm}]
kk=2; [mdd,ndd]=size(dd);
while ~isempty(V)
    [tmpd,j]=min(W(i,V)); tmpj=V(j);
    for k=2:ndd
        [tmp1,jj]=min(dd(1,k)+W(dd(2,k),V));
        tmp2=V(jj); tt(k-1,:)=[tmp1,tmp2,jj];
    end
    tmp=[tmpd,tmpj,j;tt]; [tmp3,tmp4]=min(tmp(:,1));
    if tmp3==tmpd
        ss(1:2,kk)=[i;tmp(tmp4,2)];
    else
        tmp5=find(ss(:,tmp4)~=0); tmp6=length(tmp5);
        if dd(2,tmp4)==ss(tmp6,tmp4)
            ss(1:tmp6+1,kk)=[ss(tmp5,tmp4);tmp(tmp4,2)];
        else
            ss(1:3,kk)=[i;dd(2,tmp4);tmp(tmp4,2)];
        end
    end
    dd=[dd,[tmp3;tmp(tmp4,2)]]; V(tmp(tmp4,3))=[];
    [mdd,ndd]=size(dd); kk=kk+1;
end
S=ss; D=dd(1,:);
\end{lstlisting}

\end{document}